\documentclass[11pt,a4paper]{article}

% Preambulo

\usepackage[utf8]{inputenc}
\usepackage[T1]{fontenc}
\usepackage[spanish]{babel}
\usepackage{csquotes}
\usepackage[margin=2.5cm]{geometry}
\usepackage{graphicx}
\usepackage{booktabs}
\usepackage{siunitx}
\usepackage{amsmath}
\usepackage{amssymb}
\usepackage{xcolor}
\usepackage{listings}
\usepackage{caption}
\usepackage{longtable}
\usepackage{array}

\renewcommand{\refname}{Bibliograf\'ia}

\usepackage{hyperref}
\hypersetup{
  colorlinks=true,
  linkcolor=blue!50!black,
  citecolor=blue!50!black,
  urlcolor=blue!50!black,
  pdftitle={PE-U4: pandas vs PySpark y comprobacion de la Ley de Amdahl (ACC)},
  pdfauthor={Cristhian Pacheco, Robinson Cando, Ernesto Luna}
}

% --- Estilo de listados de codigo (Python) ---
\lstdefinestyle{pythonstyle}{
  language=Python,
  basicstyle=\ttfamily\footnotesize,
  keywordstyle=\color{blue!60!black}\bfseries,
  commentstyle=\color{gray!70!black}\itshape,
  stringstyle=\color{teal!70!black},
  numbers=left,
  numberstyle=\tiny\color{gray},
  stepnumber=1,
  numbersep=6pt,
  breaklines=true,
  breakatwhitespace=false,
  showstringspaces=false,
  frame=single,
  framerule=0.3pt,
  captionpos=b,
  tabsize=4,
}
\lstset{style=pythonstyle}

% --- Comando para marcar campos institucionales que el equipo debe completar ---
\newcommand{\completar}[1]{\colorbox{yellow!60}{\textbf{[COMPLETAR: #1]}}}

\title{PE-U4: pandas vs PySpark y comprobacion experimental de la Ley de Amdahl}
\author{}
\date{}

\begin{document}

% PORTADA

\begin{titlepage}
\centering
{\large UNIVERSIDAD T\'ECNICA ESTATAL DE QUEVEDO\par}
{\large Facultad de Ciencias de la Computaci\'on\par}
{\large Carrera de Ingenier\'ia de Software\par}

\vspace{1.2cm}
{\Large\bfseries PE-U4: Procesamiento Distribuido con Apache Spark\\[2pt]
Comprobaci\'on Experimental de la Ley de Amdahl\\[2pt]
sobre un Pipeline de Tickets de Soporte T\'ecnico ISP (ACC)\par}

\vspace{0.8cm}
\begin{tabular}{@{}ll@{}}
\textbf{C\'odigo de actividad:} & GA-SUM-05 / PE-U4 \\
\textbf{Asignatura:} & Aplicaciones Distribuidas (ISR-701) \\
\textbf{Unidad:} & Unidad 4 \\
\textbf{Nivel / Paralelo:} & S\'eptimo de Ingenier\'ia de Software ``A'' \\
\textbf{Docente:} & Gleiston C. Guerrero-Ulloa, M.Sc. \\
\textbf{Per\'iodo acad\'emico:} & 2026--2027 PPA \\
\end{tabular}

\vspace{0.8cm}
\textbf{PFC de referencia:} ACC --- Sistema de Gesti\'on de Soporte T\'ecnico ISP

\vspace{0.6cm}
{\large\textbf{Integrantes}\par}
\vspace{0.2cm}
\resizebox{\textwidth}{!}{%
\begin{tabular}{@{}lll@{}}
\toprule
\textbf{Nombre} & \textbf{PFC de origen} & \textbf{Rol asumido en esta pr\'actica} \\
\midrule
Cristhian Daniel Pacheco C\'ardenas & ACC & Implementaci\'on PySpark, protocolo de medici\'on y escalado de T3 \\
Robinson Rodrigo Cando Moreno & ACC & An\'alisis de Amdahl, equivalencia y umbral de rentabilidad \\
Ernesto Gregory Luna Mora & BCEL & Repositorio, reproducibilidad, figuras y verificaci\'on final \\
\bottomrule
\end{tabular}%
}

\vspace{0.8cm}
\begin{minipage}{0.92\textwidth}
\textbf{Justificaci\'on t\'ecnica.} El dominio ACC (Sistema de Gesti\'on de Soporte
T\'ecnico ISP) opera sobre un flujo continuo de tickets de soporte --cortes de
servicio, reclamos de velocidad, fallas de equipo, solicitudes de instalaci\'on
y consultas de facturaci\'on-- generado las 24 horas por una base de clientes
activa. Incluso un proveedor de tama\~no medio, con algunos miles de clientes,
acumula en pocos meses de operaci\'on cientos de miles de registros de
tickets, cada uno con m\'ultiples atributos temporales, categ\'oricos y de
estado. A esa escala, el an\'alisis explorat\'orio, los reportes gerenciales
recurrentes y las agregaciones por zona o categor\'ia dejan de ser viables con
herramientas de una sola m\'aquina: el tiempo de c\'omputo crece de forma que
compromete la oportunidad de la decisi\'on operativa. Este trabajo usa
exactamente ese volumen (600\,000 tickets sint\'eticos, semilla fija) para
comprobar, con datos propios, si un motor distribuido como Apache Spark ofrece
una ventaja medible frente a pandas. La decisi\'on de negocio concreta que se
apoyar\'ia en este pipeline es identificar qu\'e categor\'ias de incidencia
concentran m\'as tickets cr\'iticos por zona operativa, para decidir d\'onde
reforzar personal t\'ecnico de campo o priorizar mantenimiento de
infraestructura de red --una decisi\'on operativa medible, no una mejora
gen\'erica del servicio.
\end{minipage}

\vfill
{\small Repositorio: \url{https://github.com/CristhianP03/pe-u4-spark-equipoA}\par}
\end{titlepage}

\tableofcontents
\newpage

\section*{Resumen}
\addcontentsline{toc}{section}{Resumen}

La Ley de Amdahl (1967) establece el l\'imite te\'orico de aceleraci\'on que
puede obtenerse al paralelizar un programa cuya fracci\'on serial es fija, y
sigue siendo el marco de referencia para evaluar si motores modernos de
procesamiento distribuido, como Apache Spark, cumplen su promesa de
escalabilidad sobre cargas de trabajo reales. Este informe comprueba
experimentalmente dicha ley aplicando un pipeline de cinco transformaciones
(filtrado, agrupaci\'on, join, generaci\'on de variables derivadas y top-N)
sobre un dataset sint\'etico determinista de 600\,000 tickets de soporte
t\'ecnico ISP del dominio ACC, implementado en paralelo en pandas y en
PySpark~3.5.3, bajo un protocolo de medici\'on riguroso (calentamiento
descartado, cinco repeticiones cronometradas con \texttt{time.perf\_counter},
estad\'istico mediana) y con escalado de la transformaci\'on de join (T3) en
$N=1,2,4$ executors. Los resultados muestran una aceleraci\'on de PySpark
frente a pandas de entre 1.43$\times$ y 2.56$\times$ seg\'un la
transformaci\'on, con equivalencia num\'erica exacta entre ambos motores en
las cinco transformaciones (cardinalidad, columnas, sumas, medias y claves de
agrupaci\'on). El an\'alisis de Amdahl sobre T3 estima una fracci\'on
paralelizable observada de apenas 17.3\,\%, lo que fija un techo te\'orico de
aceleraci\'on de 1.21$\times$ y explica por qu\'e el rendimiento se satura ya
con $N=2$ executors --un hallazgo consistente con el \emph{shuffle} del join
como cuello de botella dominante, m\'as que con una limitaci\'on
algor\'itmica intr\'inseca. Estos resultados tienen una implicaci\'on directa
para ACC: la migraci\'on a Spark se justifica por el volumen de datos que hay
que procesar, no por el paralelismo bruto de hardware adicional.

\newpage

\section{Introducci\'on}

El crecimiento sostenido del volumen de datos operativos en sistemas
empresariales ha convertido el procesamiento distribuido en una necesidad
pr\'actica m\'as que en una opci\'on de dise\~no: cuando un conjunto de datos
supera lo que una sola m\'aquina puede procesar en un tiempo aceptable, la
\'unica v\'ia para mantener la oportunidad de las decisiones basadas en esos
datos es repartir el c\'omputo entre varias unidades de procesamiento
\cite{dean2008,zaharia2016}.

Sin embargo, repartir el trabajo entre m\'as procesadores no garantiza una
aceleraci\'on proporcional. La Ley de Amdahl, formulada en 1967, es el marco
cl\'asico que cuantifica ese l\'imite: toda tarea computacional tiene, en la
pr\'actica, una fracci\'on que no puede paralelizarse --ya sea por
dependencias de datos, coordinaci\'on entre nodos o simple l\'ogica
secuencial-- y esa fracci\'on fija un techo de aceleraci\'on que ning\'un
n\'umero adicional de procesadores puede superar. M\'as de cinco d\'ecadas
despu\'es de su formulaci\'on, la vigencia de esta ley sigue siendo objeto de
verificaci\'on emp\'irica cuando se aplica a motores modernos de
procesamiento distribuido como Apache Spark, cuyo modelo de ejecuci\'on en
memoria y evaluaci\'on perezosa difiere sustancialmente de los sistemas para
los que Amdahl escribi\'o originalmente su argumento.

El objetivo de este trabajo es comprobar experimentalmente la Ley de Amdahl
--y su complemento, la Ley de Gustafson-Barsis-- aplicando un pipeline real
de cinco transformaciones de datos sobre un dataset sint\'etico de 600\,000
tickets de soporte t\'ecnico ISP, correspondiente al dominio ACC (Sistema de
Gesti\'on de Soporte T\'ecnico ISP), implementado en paralelo en pandas
(l\'inea base de una sola m\'aquina) y en PySpark (motor distribuido), con
escalado controlado de una de las transformaciones en $N=1,2,4$ executors.

El documento se organiza de la siguiente manera: la Secci\'on~\ref{sec:fundamento}
desarrolla el fundamento te\'orico del paralelismo, MapReduce/Hadoop, la
arquitectura de Spark y los modelos de computaci\'on en la nube; la
Secci\'on~\ref{sec:procedimiento} documenta el procedimiento aplicado
(dataset, configuraci\'on de Spark, implementaci\'on y verificaci\'on de
equivalencia); la Secci\'on~\ref{sec:resultados} presenta los resultados
cuantitativos; la Secci\'on~\ref{sec:analisis} desarrolla el an\'alisis
cr\'itico de la fracci\'on serial observada y el umbral de rentabilidad; la
Secci\'on~\ref{sec:preguntas} responde las preguntas de an\'alisis del Anexo
A; y el documento cierra con las conclusiones individuales del equipo y la
declaraci\'on de uso de inteligencia artificial generativa.
